\chapter{Architecture}
\label{chap-architecture}

\section{Une seule règle : les dépendances pointent vers l'intérieur}

L'architecture tient en une phrase : \textbf{le moteur de jeu ne sait pas comment il est
    affiché ; les interfaces connaissent le moteur ; rien ne connaît les interfaces}
(\cref{fig-couches}).

\begin{figure}[ht]
    \centering
    \begin{tikzpicture}[node distance=7mm and 9mm]
        \node[coeur, minimum width=11.2cm, minimum height=1.5cm] (core)
        {\textbf{battleship\_core} --- moteur de jeu (le \emph{modèle})\\[2pt]
            \footnotesize C11, bibliothèque standard seule, aucune allocation dynamique};

        \node[couche, above=14mm of core.north west, anchor=south west, minimum width=3.4cm,
            minimum height=1.6cm] (gui)
        {\textbf{battleship}\\[1pt]\footnotesize application GTK~3\\\footnotesize vue + contrôleur};
        \node[couche, above=14mm of core.north, anchor=south, minimum width=3.4cm,
            minimum height=1.6cm] (cli)
        {\textbf{battleship-cli}\\[1pt]\footnotesize jeu en terminal\\\footnotesize simulateur};
        \node[couche, above=14mm of core.north east, anchor=south east, minimum width=3.4cm,
            minimum height=1.6cm] (tests)
        {\textbf{tests}\\[1pt]\footnotesize unitaires\\\footnotesize et de propriétés};

        \node[externe, above=9mm of gui, minimum width=3.4cm] (gtk) {GTK~3, GLib\\[-1pt]
            \footnotesize GStreamer (facultatif)};

        \draw[fleche] (gui.south) -- (gui.south |- core.north);
        \draw[fleche] (cli.south) -- (core.north);
        \draw[fleche] (tests.south) -- (tests.south |- core.north);
        \draw[fleche, draw=ardoise, dashed] (gui.north) -- (gtk.south);

        \node[note, right=5mm of gtk, text width=6.4cm, align=left]
        {Une flèche se lit \enquote{dépend de}.\\Aucune ne part du moteur.};
    \end{tikzpicture}
    \caption{Les couches du programme. Trois exécutables sont construits sur le même moteur.}
    \label{fig-couches}
\end{figure}

Cette règle n'est pas une convention que l'on se promet de respecter : elle est
\emph{imposée par la construction}. Le moteur est une bibliothèque statique dont le seul
chemin d'inclusion est le dossier \fichier{include/}. La chaîne d'intégration continue le
compile et le teste dans un conteneur où GTK n'est pas installé (\cref{sec-ci}) : si quelqu'un
ajoutait au moteur un appel à GTK, cette étape cesserait de compiler.

Trois programmes s'appuient sur le moteur. En avoir plus d'un n'est pas un luxe. Tant qu'il
n'existe qu'une interface, rien ne prouve que le modèle en soit indépendant ; la seconde
interface en est la démonstration.

\section{Du cours au projet : modèle, vue, contrôleur}

Le découpage \emph{modèle / vue / contrôleur} est né à la fin des années 1970 dans
l'environnement Smalltalk~\cite{reenskaug1979, krasner1988}. Il sépare ce qu'un programme
\emph{sait}, ce qu'il \emph{montre} et la façon dont il \emph{réagit}.

Le cours décrit le fonctionnement d'une application graphique en trois temps~\cite{toulouse2023} :
l'utilisateur agit sur la vue, l'action correspondante se déroule dans le modèle, puis la vue
est mise à jour. Il en tire trois bibliothèques. Le \cref{tab-mvc} montre comment le
projet les réalise.

\begin{table}[ht]
    \centering
    \begin{tabularx}{\textwidth}{@{}l>{\raggedright\arraybackslash}X>{\raggedright\arraybackslash}X@{}}
        \toprule
                            & \textbf{Rôle (cours)}                                                  & \textbf{Réalisation (projet)} \\
        \midrule
        \textbf{Modèle}     & Types utilisateur et fonctions de manipulation : le jeu et ses règles. &
        \fichier{include/battleship/}, \fichier{src/core/} --- plateau, tirs, déductions,
        tour de jeu, adversaire.                                                                                                     \\
        \addlinespace
        \textbf{Vue}        & Composants graphiques à organiser dans des fenêtres.                   &
        \fichier{src/ui/view.c} --- tous les \emph{widgets}, aucune règle ;
        \fichier{assets/style.css} pour l'apparence.                                                                                 \\
        \addlinespace
        \textbf{Contrôleur} & Chef d'orchestre : associe les actions aux événements,
        synchronise la vue. &
        \fichier{src/ui/controller.c} --- demande au modèle de jouer, puis resynchronise
        la vue.                                                                                                                      \\
        \bottomrule
    \end{tabularx}
    \caption{Correspondance entre l'organisation du cours et celle du projet.}
    \label{tab-mvc}
\end{table}

\subsection{Qui possède quoi ?}

Le cours discute quatre façons de regrouper le modèle et la vue dans la structure du
contrôleur, et retient celle où le modèle est \emph{référencé} et la vue \emph{possédée} :
le modèle a une existence propre, antérieure et postérieure à toute interface graphique,
alors que la vue n'a de raison d'être qu'avec son contrôleur. Le projet fait exactement ce
choix (\cref{lst-controleur})\footnote{Dans les extraits de la version~2, les commentaires,
    rédigés en anglais dans le code, sont traduits ; certains sont abrégés.}.

\begin{lstlisting}[caption={La structure du contrôleur (\fichier{src/ui/controller.c}).},
                   label={lst-controleur}]
struct bs_controller {
    GtkApplication *application;
    bs_game *model; /* emprunté : appartient à l'application */
    bs_view *view;  /* possédé  : construit et détruit par le contrôleur */
    /* ... */
};
\end{lstlisting}

Le modèle, lui, est déclaré par l'application (\fichier{src/ui/app.c}), en dehors du
contrôleur. Il pourrait être sauvegardé, rejoué ou confié à une autre interface sans que
l'interface graphique en sache rien.

\subsection{La variante \enquote{à pointeurs de fonctions}}
\label{sec-pointeurs}

Dans le corrigé du cours, le contrôleur va chercher les boutons dans la vue
(\code{vue\_tresor\_get\_boite()}) et leur attache lui-même les fonctions de rappel. Une note
ajoute qu'une autre voie, plus élégante, existe au prix de pointeurs de fonctions. C'est
cette variante que le projet met en œuvre.

La vue connecte elle-même ses signaux et traduit chaque geste de l'utilisateur en une
\emph{intention}, exprimée dans le vocabulaire du jeu, qu'elle transmet par une table de
fonctions de rappel (\cref{lst-actions}).

\begin{lstlisting}[caption={Les intentions rapportées par la vue (\fichier{src/ui/view.h}).},
                   label={lst-actions}]
typedef struct bs_view_actions {
    void (*level_chosen)(bs_difficulty level, gpointer user_data);
    void (*cell_activated)(bs_coord cell, gpointer user_data);
    void (*reshuffle_requested)(gpointer user_data);
    void (*restart_requested)(gpointer user_data);
    void (*quit_requested)(gpointer user_data);
    void (*game_over_answered)(bs_view_choice choice, gpointer user_data);
    void (*closed)(gpointer user_data);
} bs_view_actions;
\end{lstlisting}

La conséquence est forte : \textbf{aucun \code{GtkWidget} ne franchit l'interface de la
    vue}. Le contrôleur ne manipule pas de boutons ; il dit \emph{quoi} montrer --- \enquote{cette
    case est un coup au but} --- jamais \emph{comment}. Il pourrait piloter une tout autre
bibliothèque graphique sans changer d'une ligne.

\section{Les modules}

\begin{figure}[ht]
    \centering
    \begin{tikzpicture}[node distance=6mm]
        \node[boite] (types) {\textbf{types}\\[-1pt]\footnotesize coordonnées,\\[-2pt]\footnotesize statuts};
        \node[boite, right=of types] (rng) {\textbf{rng}\\[-1pt]\footnotesize hasard\\[-2pt]\footnotesize reproductible};
        \node[boite, right=of rng] (board) {\textbf{board}\\[-1pt]\footnotesize flotte, tirs,\\[-2pt]\footnotesize déductions};
        \node[boite, right=of board] (ai) {\textbf{ai}\\[-1pt]\footnotesize choix\\[-2pt]\footnotesize d'une cible};
        \node[boite, right=of ai] (game) {\textbf{game}\\[-1pt]\footnotesize tour de jeu,\\[-2pt]\footnotesize score, fin};
        \draw[fleche] (rng) -- (types);
        \draw[fleche] (board) -- (rng);
        \draw[fleche] (ai) -- (board);
        \draw[fleche] (game) -- (ai);
    \end{tikzpicture}
    \caption{Les modules du moteur forment une chaîne : chacun ne s'appuie que sur les
        précédents.}
    \label{fig-modules}
\end{figure}

À l'intérieur du moteur, les cinq modules forment une chaîne stricte (\cref{fig-modules}).
Chacun suit le patron du cours : un fichier d'en-tête qui fixe les responsabilités ---
structure de données, mode d'emploi, pré-conditions --- et un fichier source qui les réalise. Le \cref{tab-arborescence} donne l'organisation complète du dépôt.

\begin{table}[ht]
    \centering
    \small
    \begin{tabularx}{\textwidth}{@{}l>{\raggedright\arraybackslash}Xr@{}}
        \toprule
        \textbf{Dossier}                                   & \textbf{Contenu}                                           & \textbf{Lignes} \\
        \midrule
        \fichier{include/battleship/}, \fichier{src/core/} & Moteur de jeu : interface publique et réalisation          & \num{1343}      \\
        \fichier{src/cli/}                                 & Interface terminal et simulateur                           & \num{645}       \\
        \fichier{src/ui/}                                  & Application GTK : vue, contrôleur, coquille, son           & \num{1670}      \\
        \fichier{assets/}                                  & Images, sons, feuille de style (incorporés à l'exécutable) & \num{199}       \\
        \fichier{tests/}                                   & Tests unitaires, de propriétés et de l'interface terminal  & \num{1582}      \\
        \fichier{docs/}                                    & Guide d'architecture, décisions, présent rapport           &                 \\
        \fichier{cmake/}, \fichier{packaging/}             & Modules de construction, intégration au bureau             &                 \\
        \fichier{.github/}                                 & Intégration et livraison continues                         &                 \\
        \bottomrule
    \end{tabularx}
    \caption{Organisation du dépôt. Les lignes comptent le code et ses commentaires ;
        pour \fichier{assets/}, la feuille de style seule.}
    \label{tab-arborescence}
\end{table}

\begin{retenir}
    Le code des tests (\num{1582} lignes) est plus volumineux que celui du moteur qu'il vérifie
    (\num{1343} lignes, en-têtes compris). Ce rapport est habituel pour un composant dont on
    veut pouvoir garantir le comportement.
\end{retenir}
